\subsection{Der uniforme Sewing-Test wird zu einem exakt zertifizierten Auswahlsatz}
Innerhalb der unitären affinen E8-Vakuumklasse gilt
\begin{equation}
c(k)=\frac{248k}{k+30},\qquad
d_2(k)=31124\ (k\ge2),\qquad d_2(1)=4124.
\end{equation}
Die erste affine Integrabilitäts-Nullrelation tritt bei Grad $k+1$ auf; deshalb ist der Grad-zwei-Raum für $k\ge2$ noch
$248+\dim\operatorname{Sym}^2(248)=31124$.
Für die erklärte Vakuumblock-Sewinggröße
\begin{equation}
\mathcal R_\chi(k)=2^{-3c(k)/8}\chi_{k,0}(i)
=e^{c(k)\varepsilon}\sum_{n\ge0}d_n e^{-2\pi n},
\quad
\varepsilon=\frac{\pi}{12}-\frac38\log2
\end{equation}
folgt wegen $c(k)\ge31/2$ und positiver weiterer Koeffizienten die im Anhang angegebene Schranke
\begin{equation}
\mathcal R_\chi(k)\ge
e^{(31/2)\varepsilon}
\bigl(1+248e^{-2\pi}+31124e^{-4\pi}\bigr)
=1.61786701766\ldots>\frac32.
\end{equation}
Eine gröbere, vollständig rationale Zertifizierung genügt bereits:
$512<e^{2\pi}<540$ impliziert $\varepsilon>0$ und damit
\begin{equation}\label{eq:rationalsewing}
\boxed{\mathcal R_\chi(k)>
1+\frac{248}{540}+\frac{31124}{540^2}
=\frac{114161}{72900}
>\frac32,\qquad k\ge2.}
\end{equation}
Die Exponentialgrenzen werden mit den klassischen rationalen Schranken
$157/50<\pi<22/7$ und positiven Taylorreihen samt geometrisch beschränktem Rest geprüft. Der Abstand $4811/72900$ ist damit kein Rundungsbefund.

Für $k\ge2$ ist der einzelne Vakuumcharakter allerdings noch nicht die vollständige physische Replica-Antwort einer nicht-holomorphen Theorie.
Eine saubere zusätzliche Vergleichsklasse sind volle unitäre linke/rechte E8$_k$-Theorien mit nichtnegativen Sektormultiplizitäten und Vakuummultiplizität eins. Bei $\tau=i$ gilt
$Z(i)\ge\chi_{k,0}(i)^2$.
Die vollständige Zweintervallformel~\cite{headrickreplica} liefert daher
\begin{equation}
R_{2,\mathrm{voll}}(1/2)
=2^{-3c(k)/4}Z(i)
\ge\mathcal R_\chi(k)^2
>\left(\frac{114161}{72900}\right)^2
>\frac94.
\end{equation}
Eine \emph{unabhängig} hergeleitete vollständige Antwort $9/4$ würde alle höheren Level auch in dieser vollen Vergleichsklasse ausschließen. Damit wird der chiralen TFPT-Quelle kein rechter Sektor hinzugefügt. Beide Selektoren bleiben bedingt darauf, dass der Rohprozess die jeweilige Antwort und ihre Sewingbedeutung unabhängig festlegt. Ein aus E8$_1$ berechnetes $3/2$ oder $9/4$ darf E8$_1$ nicht selbst begründen.

\subsection{Der verbundene nächste Herkunftsschritt}
Die neuen Rechnungen verbinden jetzt ausdrücklich vier Ebenen:
Die gleiche Viermarkengeometrie liefert die Paarungsgewichte; diese liefern einen Viertelphasen-Anker; die markierte Quelle trägt dessen kohärente Wirkung; eine konkret definierte Auslese erzeugt die alte Clockmatrix. Derselbe vollständige Quellenraum enthält außerdem die Ladungsfasern und das auf geschlossenen projizierten Wegen verbleibende Gedächtnis.

Die Auswahlfrage ist damit enger geworden. Der ursprüngliche Prozess muss die kohärente Familien-Auslesebasis, den tatsächlichen Umgang mit Zwischenrecords und eine Bewegung beziehungsweise lokale Ordnung der vollen Fasern festlegen. Der Clockcharakter allein bewegt die Fasern nicht. Die Faserkoordinate allein liefert keine lokalen Labore. Die zusätzliche Sewing-Schranke ist ein scharfer Vergleichstest, sobald eine unabhängige Rohquellenantwort vorliegt.

\begin{center}
\begin{tikzpicture}[node distance=7mm]
\node[box,text width=13.2cm] (s) {\textbf{Eine vollständige geladene Quelle}\\Feldoperationen, Zustand, Ladungen, Phasen und Nachfahrenturm};
\node[box,below=10mm of s,xshift=-3.6cm,text width=6cm] (c) {\textbf{Clockansicht}\\Markierter Phasenschritt und gewählte Auslese ergeben $B(x)$. Kohärenzen bestimmen die spätere Rückwirkung.};
\node[box,below=10mm of s,xshift=3.6cm,text width=6cm] (f) {\textbf{Faseransicht}\\Die Dreierkoordinate ordnet Ladungen. Innere Faser und Moden unterscheiden auch geschlossene projizierte Wege.};
\node[box,below=45mm of s,text width=13.2cm] (u) {\textbf{Universalraumregel für beide Ansichten}\\Behalte alle Unterschiede, die eine zulässige spätere Operation wieder sichtbar machen kann.};
\draw[arr] (s)--(c);\draw[arr] (s)--(f);
\draw[arr] (c)--(u);\draw[arr] (f)--(u);
\end{tikzpicture}
\end{center}

\key{\textbf{Was der Gesamtlösung jetzt konkret näherkommt:}
Die vorher getrennten Clockkonstruktionen sind für alle $x$ verbunden; der Quellen-Intertwiner, die orthogonale Erweiterung der Trine und die Rückwirkung der ausgelassenen Kohärenzen sind ausdrücklich berechnet. Der volle Quellenzustand erklärt zugleich, warum ein geschlossener projizierter Weg Information behalten kann. Die noch zu beweisende Herkunft betrifft die gemeinsame Auswahl dieses Instruments und einer physischen lokalen Dynamik auf derselben Quelle. Eine vollständige Raumzeit- und Wechselwirkungsherleitung ist durch die neuen Darstellungen allein weiterhin nicht geliefert.}

\section{Gemeinsame Zeit, wirkliche Übertragung und die Auswahl des Quellprozesses}\label{sec:sharedtime}

\key{\textbf{Der neue Anhang verbindet die richtige Frage mit einer noch zu engen Ereignisklasse.}
Die gemeinsame Clockhierarchie lässt sich exakt ausführen. Ihre kontinuierliche Markovfortsetzung ist jetzt für die ganze erklärte Klasse entschieden. Gleichzeitig zeigt die tatsächliche Quellenalgebra, warum diese Hierarchie allein keine wechselwirkende Welt beschreibt und warum der vorgeschlagene räumliche Lauf nicht unverändert auf die C-Felder übertragen werden kann. Ein positiver Anschluss ist bereits in derselben Quelle vorhanden: Lokale Stromimpulse besitzen unter ihrer bestehenden konformen Zeit eine exakt berechenbare transportierte Wirkung.}

Geprüft wird der eingefügte Text \file{f548f0d4-2ae2-4067-8f27-3b19b7c82ae4}. Seine verlinkten Dateien unter \file{sandbox:/mnt/data/} wurden nicht mitgeliefert. Die genannten 37 neuen und 75 älteren Bedingungen sind deshalb keine hier übernommenen Ausführungsnachweise. Die zentralen Formeln wurden eigenständig aus den vorhandenen Quellenoperatoren rekonstruiert. Die Einzelregister-GKLS-Fortsetzung und die freie Richtung $q$ waren bereits im Rechenkern enthalten; neu geprüft werden insbesondere die gemeinsame Registerhierarchie, ihr vollständiger Einbettungsbereich und die quelltreue räumliche Übertragung.

\subsection{Der ganze Quantenkanal und die gemeinsame Ereignisregel}
Auf der tatsächlichen markierten Zweierebene gilt in der Pauli-Reihenfolge $X,Y,Z$
\begin{equation}
\Phi_q(X)=\tfrac23X,\qquad \Phi_q(Y)=qY,\qquad \Phi_q(Z)=\tfrac13Z.
\end{equation}
Die ursprüngliche Trine erfüllt $\Phi_q^*(E_i)=\sum_jB_{ij}E_j$.
Am historischen, weiterhin bedingt ausgewählten Punkt $q=1/3$ ist
\begin{equation}
\Phi_{1/3}(\rho)=\tfrac7{12}\rho+\tfrac14X\rho X
+\tfrac1{12}Y\rho Y+\tfrac1{12}Z\rho Z.
\end{equation}
Der partiell transponierte normierte Choioperator hat den Eigenwert $-1/12$. Dieser Kanal kann also Verschränkung mit einer Referenz erhalten. Die dreidimensionale vollständige Dephasierung aus Abschnitt~\ref{sec:anchorfibre} ist dagegen ein anderer Quantenprozess. Die Übereinstimmung der Matrix $B$ bleibt richtig, identifiziert aber nicht die beiden vollständigen Instrumente.

Die angegebene Einzelregister-Fortsetzung stimmt:
\begin{equation}
\mathcal L_1=\tfrac{\log6}{4}(\operatorname{Ad}_X-\Id)
+\tfrac{\log(3/2)}4(\operatorname{Ad}_Y-\Id)
+\tfrac{\log(3/2)}4(\operatorname{Ad}_Z-\Id),\qquad
 e^{\mathcal L_1}=\Phi_{1/3}.
\end{equation}
Die Operatoren liegen in der von den nativen Reflexionen erzeugten Algebra. Die Definition einer Rate setzt jedoch eine Zeiteinheit voraus und macht eine solche kohärente Kombination noch nicht zu einem aus P1/P2 ausgewählten Sprungereignis.

Die sechs bereits identifizierten nativen Wörter $u_g$ besitzen auf dieser Ebene die $S_3$-Gewichte
\begin{equation}\label{eq:sharedweights}
w(q)=\left(\tfrac12+\tfrac q6,\tfrac q6,\tfrac q6,
\tfrac1{18}-\tfrac q6,\tfrac29-\tfrac q6,\tfrac29-\tfrac q6\right),
\quad 0\le q\le\tfrac13,
\end{equation}
in der Reihenfolge $e,c,c^{-1},r_{59},r_{14},r_{13}$.
Für $N$ Register und einen gemeinsamen Record gilt
\begin{equation}\label{eq:sharedhierarchy}
\Psi_{N,q}(\rho)=\sum_gw_g(q)u_g^{\otimes N}\rho(u_g^{\otimes N})^*,
\qquad
\mathcal V_{N,q}\psi=\sum_g\sqrt{w_g(q)}\ket g\otimes u_g^{\otimes N}\psi.
\end{equation}
Direkt folgen $\mathcal V^*\mathcal V=I$ und die Verträglichkeit mit jeder partiellen Spur. Bei unabhängig erneuerten Ereignissen und ohne Zwischeneingriffe ist die gesamte ganzzahlige Entwicklung geschlossen:
\begin{equation}
\begin{split}
w_e^{*m}&=(1+h+2r+2s)/6,\\
w_c^{*m}=w_{c^{-1}}^{*m}&=(1+h-r-s)/6,\\
w_{r_{59}}^{*m}&=(1-h-2r+2s)/6,\\
w_{r_{14}}^{*m}=w_{r_{13}}^{*m}&=(1-h+r-s)/6,
\end{split}\qquad
r=(2/3)^m,\quad s=(1/3)^m,\quad h=q^m.
\end{equation}
Damit gilt $\Psi_{N,q}^m=\sum_gw_g^{*m}\operatorname{Ad}_{u_g^{\otimes N}}$ für alle $N$ und $m\ge0$; bei $m=0$ gilt $w^{*0}=\delta_e$, die Punktmasse am Identitätsereignis. Eingriffe zwischen den Schritten müssen ausdrücklich eingefügt werden. Die Faltungsformel ist kein Ersatz für solche geordneten Geschichten.

\subsection{Die gemeinsame kontinuierliche Zeit ist für die ganze Klasse entschieden}
\key{\textbf{Neuer, stärkerer Satz:}
Die Hierarchie~\eqref{eq:sharedhierarchy} besitzt genau dann eine konsistente zeitunabhängige GKLS-Fortsetzung auf denselben Registern, wenn $2/9\le q\le1/4$. Die Notwendigkeit folgt bereits aus zwei Registern und setzt den hypothetischen Generator nicht vorher als klassische oder reversible $S_3$-Mischung an.}

Zunächst bestätigt sich der einfache Endpunktzeuge des Anhangs. Für $q=1/3$ ist
\begin{equation}
w*w=(10/27,5/54,5/54,1/27,11/54,11/54).
\end{equation}
Die sechs $V_g=u_g\otimes u_g$ sind unabhängig. Ihre Gram-Matrix ist
\begin{equation}
G_{gh}=\tr(V_g^*V_h)=\diag(3I_3+\mathbf1\mathbf1^T,3I_3+\mathbf1\mathbf1^T),
\quad \det G=2916.
\end{equation}
Daher ist der Choi-Rang nach einem Schritt fünf und nach zwei Schritten sechs. Alle Tensorpotenzen $u^{\otimes N}$ mit $N\ge2$ enthalten die drei irreduziblen $S_3$-Darstellungen; der gleiche Rangsprung gilt für sämtliche dieser $N$.

Für eine endlichdimensionale zeitunabhängige GKLS-Halbgruppe ist der Krausrang für $t>0$ konstant. Eine Begründung benutzt die positiven Sprungtrajektorien: Nach Abspaltung der invertiblen driftenden Linksmultiplikation wird ihr linearer Raum von Produkten der analytisch konjugierten Sprungoperatoren erzeugt. Die Auswertung auf jedem offenen Zeit-Simplex spannt denselben analytischen Funktionenraum auf; seine Dimension hängt nicht von der positiven Intervalllänge ab. Der Rangsprung schließt somit jede solche Halbgruppe am Endpunkt aus. Die Hauptquadratwurzel liefert zusätzlich den unabhängigen negativen normierten Choi-Erwartungswert
\begin{equation}
\frac{3+\sqrt3-2\sqrt6}{20}<0.
\end{equation}

\textbf{Vollständiger Beweis im offenen Bereich.}
Für $0<q<1/3$ zerfällt der Zweiregisterraum multiplikitätsfrei als
$H=H_1\oplus H_\varepsilon\oplus H_\sigma$ mit Dimensionen $1,1,2$.
Seine orthonormale Basis kann als
\begin{equation}
\tfrac1{\sqrt2}(1,0,0,1),\quad
\tfrac1{\sqrt2}(0,1,-1,0),\quad
\tfrac1{\sqrt2}(1,0,0,-1),\quad
\tfrac1{\sqrt2}(0,1,1,0)
\end{equation}
gewählt werden. Schreibe $P_1,P_\varepsilon,P_\sigma$ für die zugehörigen Projektoren. Der Fixraum von $\Psi_{2,q}$ ist
$\C P_1\oplus\C P_\varepsilon\oplus\C P_\sigma$; seine extremalen normierten Fixzustände sind $P_1,P_\varepsilon,P_\sigma/2$.

Angenommen, $\Psi_{2,q}=e^{\mathcal K}$ mit einem beliebigen GKLS-Generator.
Die Halbgruppe kommutiert mit ihrem Zeitschritt und erhält daher dessen Fixraum. Auf dem Fixzustandssimplex ist sie positiv und spurtreu und nach Zeit eins die Identität. Sie besitzt dort positive Inversen. Solche Simplex-Automorphismen sind Permutationen; Stetigkeit erzwingt die Identität. Also fixiert sie jeden $P_i$ zu jeder Zeit. \emph{Erst daraus} folgt ihre Unitalität. Der Multiplikativitätsbereich fixer Projektoren erhält nun alle neun Räume $\operatorname{Hom}(H_j,H_i)$.

Auf diesen einzelnen Räumen sind die positiven Spektren einfach:
\begin{center}
\begin{tabular}{ll}
\toprule
Block & Spektrum\\\midrule
$11,\varepsilon\varepsilon$ & $1$\\
$1\varepsilon,\varepsilon1$ & $q$\\
$1\sigma,\sigma1,\varepsilon\sigma,\sigma\varepsilon$ & $1/3,2/3$\\
$\sigma\sigma$ & $1,2/3,q,1/3$\\\bottomrule
\end{tabular}
\end{center}
Das volle 16-dimensionale Spektrum ist dagegen degeneriert.
Blockweise muss $\mathcal K$ deshalb in derselben orthonormalen Eigenbasis diagonal sein. Seine möglichen Logarithmuszweige haben identischen reellen Anteil:
\begin{equation}
\tfrac12(\mathcal K+\mathcal K^*)=\log\Psi_{2,q}.
\end{equation}
Weil die Halbgruppe unital und spurtreu ist, sind $\mathcal K$ und sein Hilbert--Schmidt-Adjungierter beide GKLS-Generatoren. Auch ihr Mittel muss einer sein.

Setze $a=\log3$, $b=\log(3/2)$ und $c=-\log q$. Der Hauptlogarithmus hat die fünf Raten
\begin{equation}
\gamma_c=\gamma_{c^{-1}}=\frac{a+b-c}{6},\quad
\gamma_{r_{59}}=\frac{c-2a+2b}{6},\quad
\gamma_{r_{14}}=\gamma_{r_{13}}=\frac{a-b+c}{6}.
\end{equation}
Nach Projektion orthogonal zu $\operatorname{vec}I$ bleiben die anderen fünf Gruppenoperatoren linear unabhängig. Der bedingte Choioperator des Generators ist die Summe ihrer Rang-eins-Formen mit Koeffizienten $\gamma_g$. Seine Positivität ist daher äquivalent zu $\gamma_g\ge0$ für jedes $g\ne e$. Das allgemeine bedingte Choi-Kriterium ist in~\cite{wolfjoint} angegeben; die vorstehende Reduktion aller Logarithmuszweige verwendet die hier nachgerechnete Sektorstruktur.

Es folgt $2(a-b)\le c\le a+b$, also $2/9\le q\le1/4$.
Umgekehrt erzeugen genau diese nichtnegativen Raten
$\mathcal K_N=\sum_{g\ne e}\gamma_g(\operatorname{Ad}_{u_g^{\otimes N}}-\Id)$
die ganze kompatible Hierarchie. Bei $q=0$ verhindert bereits Nichtinvertierbarkeit eine endliche Generator-Exponentialdarstellung; $q=1/3$ wurde oben ausgeschlossen. Damit ist die Klasse vollständig entschieden.

Das Ergebnis geht über ein zufälliges Nullgewicht hinaus: Bei $q=3/10$ sind alle sechs Gewichte positiv, beide Choi-Ränge sechs, aber
$\gamma_{r_{59}}=\log(5/6)/6<0$. Auch dort ist jede zeitunabhängige GKLS-Einbettung auf dem Zweiregisterraum ausgeschlossen.

\picturetext{Eine Uhr kann für sich allein einen glatten Sekundenzeiger besitzen. Werden zwei Uhren durch dieselben Ereignisse angetrieben, muss auch ihre gemeinsame Reaktion zwischen den angezeigten Schritten passen. Genau diese Forderung schränkt die mögliche Fortsetzung stärker ein.}

Falls die historische Antwort $FAAF=1/27$ tatsächlich unabhängig auf diese Quelle übertragen wird, verlangt sie $q=1/3$. Die drei Anforderungen \emph{diese Antwort}, \emph{die gemeinsame Ereignisregel} und \emph{eine zeitunabhängige Markovzeit auf denselben endlichen Registern} sind dann unvereinbar. Daraus folgt keine fundamentale Diskretheit der Welt: Eine größere Quelle mit Records, zeitabhängige Ausführung und die vorhandene konforme oder modulare Quellenzeit bleiben möglich.

\subsection{Die gemeinsame Ereignisregel transportiert noch keine Wirkung zwischen Registern}
Der Bellzustand $\Omega_2=(\ket{00}+\ket{11})/\sqrt2$ besitzt unter einem gemeinsamen Ereignis Rückkehrwahrscheinlichkeit eins; unter unabhängigen Ereignissen beträgt sie $5/12$. Diese im Anhang behaupteten Werte stimmen. Die Ursache ist eine geschützte gemeinsame Symmetrie.

Eine weitergehende Gleichung entscheidet aber über die Eignung als Weltmechanik. Für jede Registerteilmenge $S$ gilt
\begin{equation}\label{eq:nosupportgrowth}
\boxed{\Psi_{N,q}^*(A_S\otimes I)
=\Psi_{|S|,q}^*(A_S)\otimes I.}
\end{equation}
Dies folgt unmittelbar aus dem Produkt jedes einzelnen Krausoperators. Jede lokale Operatoralgebra bleibt in sich; das gilt auch für jede Potenz. Eine Änderung am Eingang eines anderen Registers kann daher die lokale Ausgabe nicht beeinflussen. Separierbare Eingaben bleiben separierbar, weil ihre Ausgaben Mischungen von Produktzuständen sind. Es handelt sich um lokale Operationen mit gemeinsamem klassischem Zufall, eine bekannte Klasse~\cite{gutoskijoint}.

\key{\textbf{Gemeinsame Zufälligkeit ist eine gemeinsame Ursache, aber noch keine Wechselwirkung zwischen den Registern.}
Die Hierarchie kann vorhandene Verschränkung schützen und klassische Korrelationen erzeugen. Für sich allein erzeugt sie weder Informationsübertragung zwischen den Registern noch Verschränkung aus Produktzuständen. Sie kann deshalb kein vollständiger Ersatz für den gesuchten wechselwirkenden Feldprozess sein.}

Dies bleibt wahr, wenn man solche zustandsunabhängigen gemeinsamen Ereignisse auf Kanten oder Hyperkanten verteilt und komponiert. Eine Zeichnung mit verbindenden Kanten erzeugt noch keine transportierende Operatorwirkung. Kohärent rückgekoppelte Records, echte Mehrregisteroperatoren oder rückwirkende dynamische Vermittler wären zusätzliche Operationen; sie sind nicht bereits durch~\eqref{eq:sharedhierarchy} implementiert.

\subsection{Warum weitere simultane Symmetrieprüfungen die Ereignisgewichte nicht auswählen}
Die tatsächlichen komplexen Reflexionen sind Produkte zweier reeller E8-Weylreflexionen:
\begin{equation}
R_\alpha=s_\alpha s_{J\alpha},\qquad (\alpha,J\alpha)=0.
\end{equation}
Für die drei verwendeten Strahlen wurde dies an der ursprünglichen Wurzelwirkung geprüft. Solche Gitterisometrien besitzen passende Vakuum-erhaltende Lifts zur Gitterquelle~\cite{dongnagjoint}. Der Lift samt Kokzyklus gehört zu den Daten: Eine $S_3$-Wortgleichheit auf der markierten Zweierebene ist nicht automatisch eine kanonische $S_3$-Aktion auf allen geladenen Feldern.

Für jeden gültigen einzelnen Lift $\alpha_g$ und alle wohldefinierten Vakuumwörter gilt jedoch bereits
\begin{equation}\label{eq:vacuumweightblind}
\omega\bigl(\alpha_g(A_1)\cdots\alpha_g(A_n)\bigr)
=\omega(A_1\cdots A_n).
\end{equation}
Jede normierte Mischung dieser \emph{simultanen} Antworten ist somit von $w_g$ unabhängig. Auch beliebig viele solche Ward- und Symmetrietests können die Ereignisgewichte nicht identifizieren. Dagegen können Operationen \emph{zwischen} den Einsetzungen oder Antworten präparierter angeregter Zustände die Gewichte unterscheiden. Genau in diese zweite Klasse fällt die bedingte $q/9$-Antwort.

\picturetext{Viele genaue Fotos eines symmetrischen Zahnrads verraten nicht, wie sein Motor zwischen zwei Fotos betätigt wird. Dafür muss man eine Änderung auslösen und verfolgen, wo und wann sie wieder erscheint. Die bisherigen Strukturprüfungen bleiben richtig; ihre Unempfindlichkeit gegenüber den Gewichten erklärt, warum ihre Anzahl die dynamische Auswahl nicht automatisch schließt.}

\subsection{Der vorgeschlagene Raumlauf scheitert bei direkter Übertragung auf die Quelle}
Die neue Faserenergie ist korrekt:
\begin{equation}
E_{\min}(y)=\tfrac12|y|^2+
\begin{cases}
0,&y\in\Z^3,\ \sum y\text{ gerade},\\
1/2,&y\in\Z^3,\ \sum y\text{ ungerade},\\
5/8,&y\in(\Z+1/2)^3.
\end{cases}
\end{equation}
Sie folgt durch Minimieren der ersten fünf Ladungskoordinaten unter der E8-Parität. Ein energieerhaltender homogener Ortstransport lässt sich daher nicht einfach mit beliebig großen Familienladungen und unverändertem $L_0$ identifizieren.

Auf einem zusätzlich als Ortsgitter interpretierten $A_3^*$ ist
\begin{equation}
W(k)=e^{-ik_xX/2}e^{-ik_yY/2}e^{-ik_zZ/2}
=\sum_{\epsilon\in\{\pm1\}^3}e^{-ik\cdot\epsilon/2}A_\epsilon,
\quad A_\epsilon=\tfrac18(I+\epsilon_xX)(I+\epsilon_yY)(I+\epsilon_zZ)
\end{equation}
tatsächlich unitär und besitzt die behauptete Weylgrenze. Die beiden Nullquasienergie-Kegel bei $0$ und $(\pi,\pi,-\pi)$ besitzen entgegengesetzte Orientierungen $1/8$ und $-1/8$. Zusätzlich liegen zwei inequivalente $W=-I$-Kegel bei $(2\pi,0,0)$ und $(\pi,\pi,\pi)$. Solche freien QCA-Grenzen sind etablierte Konstruktionen~\cite{darianojo}; sie wählen weder die TFPT-Wechselwirkungen noch Massen aus.

\textbf{Der entscheidende neue Test verwendet die ursprünglichen Ladungsfaktoren.}
Für die vier C-Wurzeln $\alpha_a=((-1/2)^5;f_a)$ gilt
\begin{equation}
(\alpha_a,\alpha_b)=1+\delta_{ab},\qquad
\sum_a f_a=0,\qquad \sum_a\alpha_a=((-2)^5;0)\ne0.
\end{equation}
Die vier Wurzeln und ihre Differenzen schließen auf 20 Wurzeln vom Typ $A_4$ mit Cartanrang vier. Die Projektion hat also eine tatsächliche innere Ladungsrichtung ausgeblendet; daraus folgt keine Raumzeitdimension.

Die unitären Ladungsfaktoren auf $\ell^2(E_8)$ erfüllen
\begin{equation}
U_\alpha\ket\lambda=\varepsilon(\alpha,\lambda)\ket{\lambda+\alpha},
\qquad U_{\alpha_a}U_{\alpha_b}=-U_{\alpha_b}U_{\alpha_a}\quad(a\ne b).
\end{equation}
Das ist die originale Gitter-Kokzyklusrelation~\cite{dongnagjoint}. Der Fourierbeweis für $W$ verwendet dagegen gewöhnliche kommutierende Translationen.

Ersetzt man die acht Translationen unmittelbar durch die originalen $\pm C_a$-Ladungsshifts samt den vorhandenen Feldphasen $p_\alpha$, entsteht
\begin{equation}
\widetilde W=\sum_{\alpha=\pm\alpha_a}A_{2\pi_F(\alpha)}\otimes p_\alpha U_\alpha.
\end{equation}
Die exakte Auswertung des vorhandenen Kokzyklus ergibt in $\widetilde W^*\widetilde W$ für
$\Delta=(0,0,0,0,0;0,1,1)$ den Koeffizienten
\begin{equation}\label{eq:nativehopfailure}
\boxed{\begin{pmatrix}0&0\\1/2&0\end{pmatrix}\ne0.}
\end{equation}
Die Nullverschiebung hat weiterhin Koeffizient $I$; insgesamt bleiben zwölf nichtverschwindende zusätzliche Koeffizienten. Bereits mit vollständiger Ladung und vorübergehend unterdrücktem Kokzyklus bleiben vier. Nur die gewöhnliche projizierte Gitterrechnung erfüllt alle Unitaritätsbedingungen.

Dies widerlegt genau die unmittelbare Übertragung mit diesen Richtungskoeffizienten. Die $U_\alpha$ sind dabei lediglich die unitären Ladungsfaktoren, nicht die vollständigen Vertexfelder mit ihren Nachfahren. Andere Operatoren auf der vollständigen Faser sind durch diesen Gegenfall nicht ausgeschlossen. Ein einfaches Umbenennen der C-Felder in räumliche Hops ist jedoch keine quelltreue Lösung.

\subsection{Ein positiver gemeinsamer Prozess ist bereits in der vollständigen Quelle vorhanden}
Der stärkste nächste Ansatz verwendet die vorhandene lokale E8-Quelle selbst. Betrachte einen normierten Cartan-Strom mit
\begin{equation}
[J_n,J_m]=n\delta_{n+m,0}I,\qquad [L_0,J_n]=-nJ_n.
\end{equation}
Dies sind ihre bestehenden Level-eins- und Zeitrelationen. Für reelle glatte Kreisfunktionen seien
\begin{equation}
f_n=\frac1{2\pi}\int_0^{2\pi}f(\theta)e^{-in\theta}\,d\theta,
\quad J(f)=\sum_n f_nJ_n,
\quad W(f)=e^{iJ(f)},
\quad \sigma(f,g)=\frac1{2\pi}\int f g'.
\end{equation}
Dann gilt $[J(f),J(g)]=i\sigma(f,g)I$ und unter der bereits vorhandenen Quellenzeit
\begin{equation}
\alpha_t(W(g))=W(g_t),\qquad g_t(\theta)=g(\theta-t).
\end{equation}
Die zugrunde liegende lokale Stromnetzstruktur ist Standard~\cite{tanimotojoint}; die folgende Rechnung ist ihre direkte Anwendung auf den vorhandenen Cartan-Strom.

Präpariere lokal $\Omega_\epsilon=W(\epsilon f)\Omega$, wobei $f$ in einem echten Intervall getragen wird. Aus dem zentralen Kommutator folgen ohne Näherung
\begin{equation}\label{eq:localcausalresponse}
\boxed{\omega_\epsilon(\alpha_t(W(g)))
=e^{i\epsilon\sigma(f,g_t)}\omega(W(g)),\qquad
\omega_\epsilon(J(g_t))-\omega(J(g_t))=\epsilon\sigma(f,g_t).}
\end{equation}
Die lokale Präparation und ihre spätere Ausgabe gehören zu demselben Zustand, derselben Algebra und demselben $L_0$. Sie kostet die ebenfalls bestimmte Quellenenergie
\begin{equation}
\Delta\langle L_0\rangle
=\epsilon^2\sum_{n>0}n^2|f_n|^2
=\frac{\epsilon^2}{4\pi}\int(f')^2.
\end{equation}
Die Testfunktion ist eine gewählte Eingabe an den festgelegten Quellenprozess; hier wird kein neuer Hamiltonoperator eingeführt.

Ein strikt lokaler positiver Zeuge ist ausdrücklich konstruierbar. Wähle eine nichtkonstante reelle glatte Beule $h$ mit kleinem Träger, $f=h'$ und $g(\theta)=h(\theta+t_0)$, so dass die Träger von $f$ und $g$ anfangs getrennt sind. Dann ist $\sigma(f,g)=0$, aber
\begin{equation}
\sigma(f,g_{t_0})=\frac1{2\pi}\int(h')^2>0.
\end{equation}
Die Eingriffswirkung erreicht also tatsächlich ein anderes Intervall. Dies ist chiral-konformer Transport entlang der Naht, noch keine $3+1$-Geometrie oder bereits nachgewiesene nichtabelsche Streuung.

Die lokale Operation erhält dabei den vollen reinen Vakuumzustand nicht. Damit passt sie auch zum korrekten lokalen Ausschluss des Anhangs: Sind alle beschränkten Krausoperatoren $K_a$ in einer Algebra mit separierendem $\Omega$ und gilt $\sum_aK_a\ket\Omega\bra\Omega K_a^*=\ket\Omega\bra\Omega$, so erzwingt Reinheit $K_a\Omega=c_a\Omega$ und Separiertheit $K_a=c_aI$. Die globale Projektion auf einen angeregten Clockzustand kann deshalb nicht unverändert eine solche lokale Operation sein. Der lokale Puls liefert eine tatsächliche Alternative innerhalb der Quelle und zeigt ihren Energiepreis.

\subsection{Die fundamentale Auswahlfrage ist jetzt eine gemeinsame Antwortgleichung}
\begin{center}
\begin{tikzpicture}[node distance=7mm]
\node[box,text width=12.7cm] (source) {\textbf{Dieselbe vollständige Quelle}\quad lokale Felder, Kokzyklus, Nachfahren, Vakuum und konforme Zeit};
\node[box,below=of source,text width=12.7cm] (pulse) {\textbf{Tatsächlicher Eingriff und spätere Wirkung}\quad lokale Präparation $\rightarrow$ Quellenentwicklung $\rightarrow$ messbare Antwort};
\node[box,below=of pulse,text width=12.7cm] (map) {\textbf{Gesuchte ursprüngliche Identifikation}\quad P1/P2 wählen Quelle und Instrument; eine räumliche Fortsetzung erhält dieselben gemeinsamen Antworten};
\draw[arr] (source)--(pulse);
\draw[dasharr] (pulse)--(map);
\end{tikzpicture}
\end{center}

Die bereits explizite Antwort~\eqref{eq:localcausalresponse} verhindert, dass wir die Quelle auf eine gemeinsam rauschende Anzeige reduzieren. Eine korrekte Rekodierung muss ihre geordneten Eingriffe und späteren Wirkungen erhalten, zusätzlich zu Ladungen, Statistik, höherer Antwort und Positivität. Für den Universalraum bedeutet dies: Er bewahrt die Unterschiede von \emph{beeinflussbaren zukünftigen Experimenten}, einschließlich ihres Orts- und Zeitbezugs. Eine gemeinsame Ereignisnummer oder derselbe marginale Kanal reicht dafür nicht.

Die Arbeit am fundamentalen Abschluss wird daher auf zwei gekoppelte Herkunftsgleichungen konzentriert: Der tatsächliche ursprüngliche Nahtoperator muss die erklärte Quelle samt Zustand und Instrument liefern; die räumliche Realisierung muss deren vollständige lokale Antwortstruktur erhalten und daraus die geforderten gemeinsamen physikalischen Ausgaben gewinnen. Die lokalen Impulsantworten, die gemeinsame Clockschranke und der explizite Kokzyklusgegenfall sind jetzt konkrete Bedingungen an diese Gleichungen. Sie ersetzen sie nicht durch ein weiteres frei gewähltes Weltmodell.

\key{\textbf{Was diese Fortsetzung tatsächlich entscheidet:}
Die gemeinsame zeitunabhängige Markovklasse ist vollständig klassifiziert. Eine reine gemeinsame $S_3$-Zufallsregel ist als vollständiger wechselwirkender Prozess ausgeschlossen; die naive C-Hop-Übertragung ist ebenfalls ausgeschlossen. Zugleich ist ein echter lokaler Eingriffs- und Transportprozess in der vorhandenen Quelle ausdrücklich berechnet. Die vollständige fundamentale Weltlösung ist damit nicht bewiesen. Die Suche hat aber nun einen vorhandenen positiven Quellenprozess und genaue gemeinsame Identitätsbedingungen, an denen jede vorgeschlagene Gesamtlösung bestehen muss.}


\section{Der konstruktive Rückweg: Eine lokale Feldregel bestimmt Quelle, Zustand und Zeit}\label{sec:inverseorigin}

\key{\textbf{Die entscheidende Verdichtung der Ursprungsfrage.} Die bereits gewählte, markierte E8-Level-1-Quelle lässt sich zu genau ihrer passenden Theorie mit sechzehn reellen Fermionen zurückverfolgen. Dabei werden Vakuum, vollständige Feldantworten und Konformalzeit gemeinsam rekonstruiert. Der Herkunftsnachweis muss deshalb nicht vier unabhängig gewählte Zutaten rechtfertigen. Er muss einen bestimmten lokalen Feldkeim im ursprünglichen P1/P2-Prozess nachweisen. Der folgende Satz beschreibt diesen Keim und beweist, was daraus folgt.}

\subsection{Das vorhandene E8 enthält den Rückweg bereits in seiner Markierung}
Die ursprüngliche Verklebung enthält die Zwischenstufe
\begin{equation}
D_5\oplus D_3\subset D_8\subset E_8,
\qquad D_8=\{n\in\Z^8:\textstyle\sum_i n_i\text{ gerade}\}.
\end{equation}
Sie wird nicht nachträglich ausgesucht: Im ursprünglichen $\Z_4$-Glue sind es die beiden geraden Klassen. Dessen Charakter $G$ hat auf den Klassen $j$ den Wert $i^j$; der Fixpunkt von $G^2$ ist die $D_8$-Gitteralgebra. In den ursprünglichen verdoppelten Wurzelkoordinaten $r=2\alpha$ kann man $j=\sum_{i=1}^5r_i\bmod4$ benutzen. Die vier Wurzelklassen enthalten $52,64,60,64$ Elemente. Die geraden Klassen liefern die $112$ Wurzeln $\pm e_i\pm e_j$; mit den acht Cartanrichtungen sind das die $120$ Ströme von $\mathfrak{so}(16)$. Die beiden ungeraden Klassen liefern den ursprünglichen Halbspinor mit $128$ Komponenten. Dies ist dieselbe ursprüngliche Verzweigung aus \file{tfpt_1_architecture_e8.tex}, nicht ein neuer E8-Ansatz.

Die vier Sektoren des $D_8$-Gitters sind
\begin{equation}
0,\quad v=e_1+D_8,\quad s=\tfrac12(1,\ldots,1)+D_8,\quad c=s+v,
\qquad (h_0,h_v,h_s,h_c)=(0,\tfrac12,1,1).
\end{equation}
Hier bezeichnet $h$ das niedrigste konforme Gewicht, also die halbe minimale quadrierte Gitterlänge. Genau $v$ hat den fermionischen Drehfaktor $e^{2\pi ih}=-1$. Seine Erweiterung ist ausdrücklich
\begin{equation}\label{eq:inverseZ8}
\boxed{D_8\cup(D_8+v)=\Z^8,\qquad V_{\Z^8}\cong\mathcal F_{16}.}
\end{equation}
Die zweite Gleichung ist die Boson-Fermion-Isomorphie: Jede der acht Einheitsrichtungen liefert ein komplexes Fermion, also zwei reelle Majoranafelder. Die explizite Rang-eins-Abbildung und ihre Tensorpotenzen stehen in~\cite{barronfermion}. Die sechzehn Norm-eins-Vektoren $\pm e_i$ liefern die Felder vom Gewicht $1/2$.

Die Erweiterung ist bis zu Feldphasen und markierungserhaltender Isomorphie festgelegt. Die Standard-Kokzyklen der ungeraden Gitterbeschreibung dürfen dabei auf $D_8$ durch eine Eins-Koketten-Umphasung an den ursprünglichen Kokzyklus angeglichen werden; sie dürfen nicht still an dessen Stelle treten. Für die zurückgewonnenen C/X-Antworten wird weiterhin der originale E8-Kokzyklus verwendet.

\begin{center}
\begin{tikzpicture}[node distance=8mm]
\node[box,text width=12.9cm] (m) {\textbf{Ursprünglich markiertes $E_{8,1}$}\\$D_5+A_3$, Glue und ausgewählter Halbspinor $s$};
\node[box,below=of m,text width=12.9cm] (d) {\textbf{Gemeinsame gerade Quelle $(D_8)_1$}\\Fixpunkt der halben Glue-Wirkung; derselbe Vakuum- und Stress-Tensor};
\node[box,below=11mm of d,xshift=-3.45cm,text width=5.6cm] (f) {\textbf{Fermionische Erweiterung}\\$0\oplus v$, sechzehn Majoranas\\$\mathcal H_F=\mathcal H_0\oplus\mathcal H_v$};
\node[box,below=11mm of d,xshift=3.45cm,text width=5.6cm] (e) {\textbf{Ursprüngliche bosonische Erweiterung}\\$0\oplus s$, wieder $E_{8,1}$\\$\mathcal H_{E_8}=\mathcal H_0\oplus\mathcal H_s$};
\draw[arr] (m)--(d);\draw[arr] (d)--node[left,font=\scriptsize]{einzig fermionischer Sektor}(f);\draw[arr] (d)--node[right,font=\scriptsize]{markierter Spinorsektor}(e);
\end{tikzpicture}
\end{center}

Die beiden Erweiterungen besitzen unterschiedliche Hilberträume. Auf dem ganzen E8-Raum wirkt das feste $D_8$-Unterobjekt in der Darstellung $0\oplus s$; sein eigener zyklischer Vakuumraum ist $\mathcal H_0=\overline{V_{D_8}\Omega}$. Die Majoranafelder sind keine bisher übersehenen lokalen $h=1/2$-Felder im bosonischen E8-Vakuumraum. Der markierte Halbspinor wird beim Rückweg mitgeführt und führt über $D_8+s$ wieder zur ursprünglichen E8-Quelle~\cite{chuzheng}. Insbesondere gehören die ursprünglichen C/X-Felder zum Spinor-Erweiterungssektor; sie sind nicht die neuen elementaren Majoranafelder.

